\section{Removing self-duality from the ordinary alternating shape restriction}
\label{sec:ordinary-alternating-sharp}

\paragraph{Status.}
Complete new root derivation with author whole-proof reading and an
independent whole-proof review by /root/g34\_eightfolds, both PASS
on 21 November 2026. The independent symbolic and permutation
certificate also passed. The finite checks below do not substitute
for the all-degree argument. No new finite character table is a premise.

\begin{lemma}[HD-ALT-03a: three-cycle and double-transposition commutators]
\label{HD-ALT-03a}
Let $n\ge15$ and let $W$ be a nontrivial irreducible ordinary
$A_n$-module of degree $a$. Put
\[
 u=\frac{\chi((123))}{a},\qquad
 v=\frac{\chi((12)(34))}{a}.
\]
For every nonscalar finite-order linear normalizer $T$ of the
represented core,
\[
 \nu(T)\ge a\min\left\{\frac{1-u}{3},\frac{1-v}{4}\right\}.
                                                               \tag{ALTN.1}
\]
All scalar phases are allowed. The conclusion includes both
constituents of a split Specht restriction, whether real or nonreal.
\end{lemma}

\begin{proof}
First suppose $W$ is the irreducible restriction of a Specht
module $S^\lambda$ of $S_n$. Every normalizer is scalar times
a matrix in this symmetric extension: its projective action lies
in $\operatorname{Aut}(A_n)=S_n$, and removing the implementing
symmetric matrix leaves a scalar by Schur's lemma. Let $g\in S_n$
be the implementing permutation, which is nonidentity.

If $g$ has a cycle of length at least three, choose two consecutive
letters $b,g(b)$ in that cycle and set $x=(b\ g(b))$. Its conjugate
by $g$ shares exactly one letter with $x$, so $[g,x]$ is a
three-cycle. If all nontrivial cycles of $g$ have length two and
$g$ has a fixed letter $c$, use $x=(b\ c)$ with $g(b)\ne b$;
again the commutator is a three-cycle. The only remaining case
is a fixed-point-free involution. Choose letters $b,c$ in two
different two-cycles and set $x=(b\ c)$. Its conjugate is the
disjoint transposition $(g(b)\ g(c))$, so $[g,x]$ is a double
transposition.

For any invertible matrices $A,B$ and any scalar $z$,
\[
 \operatorname{rank}([A,B]-1)
 =\operatorname{rank}(AB-BA)
 \le2\operatorname{rank}(A-z).
                                                               \tag{ALTN.2}
\]
Indeed $AB-BA=(A-z)B-B(A-z)$. Minimize over $z$ to get a bound
by $2\nu(A)$. A three-cycle has its two nontrivial eigenvalues
with equal multiplicity, since it is conjugate to its inverse
inside $A_n$. Its actual moving rank is $2a(1-u)/3$.
The moving rank of a double transposition is $a(1-v)/2$.
Applying (ALTN.2) proves (ALTN.1) in this case.

Now let $\lambda$ be self-conjugate, so
$S^\lambda|_{A_n}=W\oplus W'$ with two distinct constituents
of equal degree. We use the standard restriction rule
\cite[\S2A, p.~604]{HD:BNOT2015}. An odd permutation exchanges
them. Thus no odd automorphism is implemented on $W$, and every
normalizer is scalar times a represented element $g\in A_n$.

Suppose first that an odd permutation centralizes $g$. Its
operator on $S^\lambda$ exchanges $W$ and $W'$ and commutes
with $g$. Hence the two complete spectra of $g$ are identical,
and the largest-eigenspace codimension on $S^\lambda$ is
exactly $2\nu(g|W)$. Both the three-cycle and double-transposition
classes have odd centralizers: for the former use two fixed
letters, and for the latter use one of its transpositions.
Their normalized characters on $S^\lambda$ therefore equal
$u,v$. The preceding symmetric-module argument, in dimension
$2a$, now gives (ALTN.1) on $W$ after dividing by two.
This comparison is of full spectra, not merely the trace of $g$.

Finally suppose $g$ has no odd centralizer. It has no even
cycle, since such a cycle itself is odd and centralizes $g$.
It also has no two cycles of the same odd length: swapping
those two cycles letter by letter is an odd centralizer.
Thus its cycle lengths are distinct and odd. Since $n\ge15$,
some cycle has length at least five. In that cycle set
\[
 b_0=b,\quad b_1=g(b),\quad b_2=g^2(b),\quad b_3=g^3(b),
 \qquad x=(b_0\ b_1\ b_2)\in A_n.
\]
Direct multiplication gives
$[g,x]=(b_1\ b_2\ b_3)(b_2\ b_1\ b_0)
=(b_0\ b_3\ b_1)$, a three-cycle. Both matrices now act on
$W$ itself, so (ALTN.2) gives the stronger first entry of
(ALTN.1). This exhausts every split-constituent case.
\end{proof}

\begin{lemma}[HD-ALT-03b: the double-transposition ratio]
\label{HD-ALT-03b}
Let $n\ge15$ and let $\lambda\vdash n$ satisfy
\[
             \max(\lambda_1,\lambda'_1)\le\lceil n/2\rceil.
                                                               \tag{ALTN.3}
\]
For either irreducible constituent of its restriction to $A_n$,
\[
 u<\frac14,\qquad v<\frac3{10},\qquad
                 \nu(T)\ge\frac7{40}a                 \tag{ALTN.4}
\]
for every nonscalar linear normalizer element.
\end{lemma}

\begin{proof}
The bound on $u$ is HD-ALT-01b. Write
\[
 t=\sum_{(i,j)\in\lambda}(j-i),\quad
 S=\sum_{(i,j)\in\lambda}(j-i)^2,\quad
 k=\lceil n/2\rceil.
\]
Let $C$ be the sum of all transpositions in $\mathbf C[S_n]$.
It is $\sum_iJ_i$ and acts on $S^\lambda$ by the scalar $t$,
by the same published content-eigenvalue theorem
\cite[Theorem~4.2; author p.~9]{HD:CMN2022} used in HD-ALT-01b.
In the group algebra,
\[
 C^2=\binom n2\,1+3K_3+2K_{22},
                                                               \tag{ALTN.5}
\]
where $K_3$ and $K_{22}$ sum all three-cycles and all double
transpositions, respectively. The coefficients count the ordered
factorizations into two transpositions. Their class sizes are
$n(n-1)(n-2)/3$ and $n(n-1)(n-2)(n-3)/8$.
Using (ALT.6) in (ALTN.5) gives the exact specialization
\[
 v=\frac{4\bigl(t^2+n(n-1)-3S\bigr)}
           {n(n-1)(n-2)(n-3)}.                            \tag{ALTN.6}
\]
The nonsplitting argument in HD-ALT-03a shows that the same ratio
holds for either constituent when the restriction splits.

The elementary row/column expression for $t$ is
\[
 t=\sum_i\binom{\lambda_i}{2}
                   -\sum_j\binom{\lambda'_j}{2}.
\]
Since every row is at most $k$, the first sum is at most
$n(k-1)/2$. For every positive integer $h$, $\binom h2\ge h-1$;
therefore the second sum is at least $n-\lambda_1\ge n-k$.
Applying the same argument to the transpose gives
\[
 |t|\le\frac{n(k-1)}2-(n-k)
                     \le\frac{(n-1)(n-2)}4.              \tag{ALTN.7}
\]
For odd $n$ the final inequality is equality of the displayed
upper bounds; for even $n$ it is strict. Dropping the nonpositive
term $-3S$ from (ALTN.6) now gives
\[
 v\le\frac{(n-1)(n-2)^2+16n}{4n(n-2)(n-3)}<\frac3{10}.
                                                               \tag{ALTN.8}
\]
For the final strict inequality, after multiplication by the
positive denominator the difference is
$2(n^3-5n^2-84n+20)$. At $n=15$ the polynomial in parentheses
is $1010$, and its forward difference is
$3n^2-7n-88>0$ for $n\ge15$.
Thus the first term in (ALTN.1) is greater than $1/4$ and the
second greater than $7/40$. The stated weak bound follows.
\end{proof}

\begin{theorem}[HD-ALT-03c: all balanced ordinary shapes are excluded]
\label{HD-ALT-03c}
Suppose an inducing block of an actual rank-nine Weil sixfold
finite-cover source has an irreducibly acting normal central-product
core of ordinary alternating groups $A_n$, with $n\ge15$,
and every factor character has a Specht label satisfying (ALTN.3).
Then these data cannot occur. The assertion includes all scalar
phases, compatible outer automorphisms and factor permutations.
It does not require essential self-duality or nonsplit polarization.

In the transitive forced-core setting D-BI-19, every remaining
ordinary alternating source with $n\ge15$ therefore has a row or
column longer than $\lceil n/2\rceil$. Its factor restriction
is irreducible and extends to a real Specht representation of $S_n$;
in particular the full inducing image is essentially self-dual.
No existence or dominance is asserted for these remaining long shapes.
\end{theorem}

\begin{proof}
Let $e$ be the total inducing degree. For a factor-preserving
nonscalar normalizer, (ALTN.4) passes to the tensor product by
fixing all other eigenvalues. A nontrivial factor permutation
has the stronger quarter bound by the weighted tensor-cycle
argument used in D-SPOR-03. Genuine scalar inertia moves all
$e$ dimensions. Thus every genuine inertia moves at least $7e/40$.

Each factor has degree at least $n^3/27\ge125$. Indeed the precise
published consequence of Rasala's classification in
\cite[proof of Theorem~9.7, p.~312]{HD:GKT2019} says that a
nonlinear Specht character of degree below $2n^3/27$ has, up to
transpose, one of the three shapes $(n-1,1),(n-2,2),(n-2,1,1)$.
None satisfies (ALTN.3). Its degree is either the factor degree
or twice that degree, proving the stated lower bound.

The actual intermediate-cover Euler identity and period
factorization give genus zero, at least twelve genuine inertias,
and total moving rank $2e+6$. Genus one has at most six branches
and period-parameter dimension at most six; genus at least two
would make $6+2e(1-h)<0$, since $e\ge125$. Consequently
\[
            12\frac{7e}{40}\le2e+6,
                         \qquad e\le60,                  \tag{ALTN.9}
\]
contrary to $e\ge125$.

For the final structural conclusion, a self-conjugate partition
has $2\lambda_1-1\le n$, so it satisfies (ALTN.3). Such partitions,
including both nonselfdual and selfdual split constituents, have
therefore been excluded. A remaining partition is not self-conjugate;
its restriction to $A_n$ is the restriction of an irreducible
real Specht representation of $S_n$. The transitive tensor factors
are thus self-dual, and the irreducible full core has a unique
invariant line of nondegenerate bilinear forms. Every normalizer
preserves that line, proving essential self-duality. This conclusion
is a necessary representation restriction, not a new geometric source.
\end{proof}

\paragraph{Verification scope.}
The short group-algebra identities, scalar-normalizer argument,
and split-constituent comparison are proved above. The retained
exact script checks the polynomial bounds, the commutator cycle
identities, and all Young-diagram content and ratio formulas in a
finite range. The all-$n$ and all-tensor-exponent statements use
the inequalities in the proof, not extrapolation. The published
content theorem, Specht restriction rule, automorphism theorem and
precise small-degree classification remain cited inputs.
